\documentclass[11pt]{article}

\usepackage[a4paper,margin=27mm]{geometry}
\usepackage[T1]{fontenc}
\usepackage{lmodern}
\usepackage{microtype}
\usepackage{amsmath,amssymb,amsthm,mathtools}
\usepackage{booktabs,longtable,array,tabularx}
\usepackage{enumitem}
\usepackage{xcolor}
\usepackage{listings}
\usepackage[hidelinks]{hyperref}
\usepackage{url}

\definecolor{leanblue}{RGB}{20,65,115}
\definecolor{codegray}{RGB}{246,247,249}
\definecolor{warningred}{RGB}{145,25,35}
\definecolor{gateblue}{RGB}{230,240,250}

\hypersetup{
  pdftitle={Guideline for an End-to-End Lean 4 Formalisation of H1 <= 236},
  pdfauthor={Project specification},
  pdfsubject={Lean 4, mathlib, bounded gaps between primes, formal verification}
}

\lstdefinelanguage{Lean}{
  morekeywords={theorem,lemma,def,structure,class,instance,namespace,end,where,
    by,have,show,exact,apply,intro,constructor,refine,obtain,let,if,then,else,
    match,with,fun,forall,exists,import,open,scoped,noncomputable,set_option,
    example,abbrev,inductive,deriving},
  sensitive=true,
  morecomment=[l]{--},
  morecomment=[s]{/-}{-/},
  morestring=[b]"
}

\lstset{
  language=Lean,
  basicstyle=\ttfamily\small,
  keywordstyle=\color{leanblue}\bfseries,
  commentstyle=\color{gray!75!black},
  stringstyle=\color{green!35!black},
  backgroundcolor=\color{codegray},
  frame=single,
  rulecolor=\color{gray!35},
  columns=fullflexible,
  keepspaces=true,
  breaklines=true,
  showstringspaces=false,
  xleftmargin=0.6em,
  xrightmargin=0.6em,
  aboveskip=0.8em,
  belowskip=0.8em
}

\newtheorem{principle}{Principle}
\newtheorem{gate}{Acceptance gate}
\newtheorem{warning}{Warning}
\newcommand{\Hone}{H_1}
\newcommand{\leanname}[1]{\texttt{#1}}
\newcommand{\pass}{\textsc{pass}}
\newcommand{\fail}{\textsc{fail}}
\newcommand{\todo}{\textsc{open}}
\newcommand{\mustpass}{\textsc{must pass}}

\title{Project Guideline for an End-to-End Lean 4 Formalisation\\
of the Target Bound \(\Hone\leq 236\)}
\author{Formalisation and release specification}
\date{Version 1.0 --- 3 September 2026}

\begin{document}

\maketitle

\begin{abstract}
This document specifies how to turn the present manuscript-conditional,
computer-assisted argument for \(\Hone\leq236\) into a genuinely
unconditional Lean 4 theorem, developed with ChatGPT/Codex in Visual Studio
Code and released with its paper on arXiv.  It fixes the formal statement,
trust boundary, repository layout, dependency graph, certificate architecture,
development protocol, audit gates, and release package.  It is a guideline and
project contract, not a claim that the bound has already been formalised.

The decisive requirement is that the final theorem must not assume the
Stadlmann sieve input, the negative-\(\omega_0\) repair, any numerical
certificate, or any analytic estimate as a project-specific axiom.  Each must
be proved in Lean from definitions and already kernel-checked library results.
The final declaration must build in a clean checkout, and its axiom audit may
contain only Lean's standard foundational axioms.  This makes the project much
larger than translating the finite certificate: the Bombieri--Vinogradov,
Maynard-sieve, Harman-minorant, and Type-I/II/III layers are the main work.
\end{abstract}

\tableofcontents

\section{Outcome and non-negotiable meaning of ``unconditional''}

The target deliverable is a Lean project whose public theorem has no hypotheses
other than definitions already appearing in its statement.  At mathematical
level it should imply
\[
  \Hone=\liminf_{n\to\infty}(p_{n+1}-p_n)\leq236.
\]
At formal level the primary theorem should use an elementary, unambiguous
``infinitely often'' formulation; a separate bridge may identify it with the
usual liminf notation.

\begin{principle}[Formalisation changes the status only at the last gate]
Until the analytic input has itself been constructed by proved Lean terms, the
correct public description remains ``conditional formalisation'' or
``formalisation in progress.''  A Lean theorem of the form
\leanname{conditional\_H1\_le\_236 (h : AnalyticInput) : ...} is useful, but it
does not make the result unconditional.
\end{principle}

For this project, ``unconditional'' means all of the following.

\begin{enumerate}[label=U\arabic*.,leftmargin=3.2em]
\item No unproved conjecture or manuscript proposition occurs as a hypothesis
  of the final theorem.
\item No project declaration made with \leanname{axiom}, no
  \leanname{sorry}/\leanname{admit}, and no externally asserted numerical fact
  occurs in the transitive dependency closure of the final theorem.
\item External solvers and interval packages may \emph{generate} data, but a
  proved Lean checker must verify every datum needed by the theorem.
\item All analytic results cited on paper, including ``standard'' endpoint and
  uniformity arguments, are either imported as already proved Lean theorems or
  proved in this project.
\item A clean build, axiom audit, and independent replay succeed with a pinned
  Lean and mathlib snapshot.
\end{enumerate}

It is legitimate to trust Lean's kernel and the pinned, kernel-checked mathlib
release.  It is also normal for \leanname{\#print axioms} to report the built-in
axioms \leanname{propext}, \leanname{Classical.choice}, and
\leanname{Quot.sound}.  Any \leanname{sorryAx}, compiler-trust computation
axiom, or project-specific axiom is a release blocker.  Lean's official proof
validation guide explains this distinction and recommends both axiom printing
and independent rechecking \cite{LeanValidation}.

\begin{warning}[Scale]
The current mathematical dossier isolates a short finite crossing, but it
imports a very deep analytic framework.  The end-to-end formalisation should be
budgeted as a multi-person-year research and engineering project unless the
initial dependency audit finds substantially more reusable formalised analytic
number theory than is presently documented.  The existing mathlib Selberg
sieve is valuable infrastructure, not a ready-made Maynard--Stadlmann theorem
\cite{MathlibSelberg}.
\end{warning}

\section{Freeze the mathematical target before coding}

\subsection{The robust primary statement}

Define consecutive primes without relying on a particular enumeration:

\begin{lstlisting}
import Mathlib

namespace PrimeGaps236

def ConsecutivePrimes (p q : Nat) : Prop :=
  Nat.Prime p /\ Nat.Prime q /\ p < q /\
    forall r : Nat, p < r -> r < q -> Not (Nat.Prime r)

def BoundedPrimeGaps (H : Nat) : Prop :=
  forall N : Nat, exists p q : Nat,
    N <= p /\ ConsecutivePrimes p q /\ q - p <= H

theorem boundedPrimeGaps_236 : BoundedPrimeGaps 236 := by
  -- Final proof term; no assumptions and no placeholders.
  ...

end PrimeGaps236
\end{lstlisting}

The ASCII logical notation above is deliberately portable in this guideline;
the source project may use Lean's Unicode notation.  This statement is the
preferred release theorem because a human auditor can read its meaning without
unfolding filter liminf machinery.

In a second file define the indexed prime gap using mathlib's
\leanname{Nat.nth Nat.Prime} and prove equivalence with the usual infinitely
often statement.  Only after that bridge is established should the paper use a
Lean theorem named, for example,
\leanname{H1\_le\_236}.  Mathlib already supplies the prime predicate, nth
prime, prime-counting definitions, and elementary connections among them
\cite{MathlibPrimeCounting}.

\subsection{Formal source of the number 236}

The last step must be factored into two independent statements:
\begin{enumerate}
\item a general theorem saying that an admissible \(48\)-tuple satisfying the
  verified sieve criterion has infinitely many translates containing two
  primes; and
\item a finite theorem that the explicit tuple
\[
\begin{split}
\mathcal H=\{&0,6,8,14,18,24,26,48,50,54,56,60,66,68,74,78,80,84,\\
&90,96,98,104,110,116,120,126,134,138,144,150,158,164,168,176,\\
&180,186,188,194,200,204,206,210,216,224,228,230,234,236\}
\end{split}
\]
is admissible, has cardinality \(48\), minimum \(0\), and maximum \(236\).
\end{enumerate}
The second theorem should be completely computational and kernel-checkable.
For primes at most \(48\), verify an omitted residue; for larger primes,
cardinality alone prevents coverage of every residue class.

\subsection{Corrected analytic proposition}

The informal source must first be rewritten as one fully quantified theorem.
The corrected theorem must make the following branch explicit:
\[
 \omega_0<0 \quad\Longrightarrow\quad
 \text{use a Bombieri--Vinogradov estimate strictly below level }1/2,
\]
whereas the four-bin Type-IIc condition is used only for
\(0\leq\omega_0\leq\omega\).  The endpoint \(\omega_0=0\) must be assigned to
exactly one branch.  It is not enough to cite the prose repair in the current
dossier: the modulus, dyadic interval, logarithmic loss, and uniform constants
must all appear in the statement proved by Lean.

\begin{gate}[Statement review]
Two analytic number theorists and one Lean reviewer should approve the
unfolded meaning of \leanname{BoundedPrimeGaps 236}, the exact corrected sieve
statement, and the quantifier order of every asymptotic parameter before the
main formalisation begins.  Record the approved statements in a tagged commit.
\end{gate}

\section{Toolchain and VS Code setup}

Lean's installation guide recommends \leanname{elan} as the toolchain manager,
VS Code with the official \leanname{leanprover.lean4} extension, and
\leanname{lake} for projects \cite{LeanInstall}.  A current mathlib project can
be created with:

\begin{lstlisting}[language=bash]
code --install-extension leanprover.lean4
lake +leanprover-community/mathlib4:lean-toolchain \
  new PrimeGaps236 math
cd PrimeGaps236
lake exe cache get
lake build
code .
\end{lstlisting}

Immediately after creation:

\begin{enumerate}
\item Commit \leanname{lean-toolchain}, \leanname{lakefile.toml} (or
  \leanname{lakefile.lean}), and \path{lake-manifest.json}.
\item Record the exact Lean version, mathlib commit hash, operating system, and
  architecture in \leanname{REPRODUCING.md}.
\item Set \leanname{set\_option autoImplicit false} in project modules, so a
  misspelled identifier does not silently create an unintended variable.
\item Turn warnings into CI failures for project code.  A use of
  \leanname{sorry} must never be merged simply because the editor shows a
  checked declaration around it.
\item Do not track \leanname{.lake/} or generated \leanname{.olean} files.
  They are caches, not the proof source.
\end{enumerate}

Do not update Lean or mathlib opportunistically.  Toolchain upgrades belong in
dedicated branches and must rebuild every module.  The release tag fixes one
manifest permanently; later upgrades produce a new release.

\section{Repository architecture}

Use a dependency-directed tree.  The following layout is a recommendation;
names may change, but the separation of logical layers should not.

\begin{lstlisting}[language={}]
PrimeGaps236/
  lean-toolchain
  lakefile.toml
  lake-manifest.json
  README.md
  REPRODUCING.md
  LICENSE
  PrimeGaps236.lean
  PrimeGaps236/
    Definitions/PrimeGaps.lean
    Foundations/FiniteSums.lean
    Foundations/Asymptotics.lean
    Foundations/ArithmeticFunctions.lean
    Foundations/ResidueClasses.lean
    Sieve/AbstractInput.lean
    Sieve/SelbergWeights.lean
    Sieve/MainTerm.lean
    Sieve/MaynardCriterion.lean
    Distribution/BombieriVinogradov.lean
    Distribution/TypeI.lean
    Distribution/TypeII.lean
    Distribution/TypeIII.lean
    Distribution/NegativeOmega.lean
    Minorant/Harman.lean
    Minorant/Mass.lean
    Support/Definitions.lean
    Support/FrontierReduction.lean
    Support/DirectCases.lean
    Support/AdaptiveCases.lean
    Certificates/RationalLP.lean
    Certificates/IntervalPolynomial.lean
    Certificates/SupportData.lean
    Certificates/IntegralData.lean
    TrialFunction/Chebyshev.lean
    TrialFunction/Integrals.lean
    TrialFunction/Crossing.lean
    Tuple/Admissible236.lean
    Final.lean
    AxiomAudit.lean
  generators/
    README.md
    export_support_certificates.py
    export_integral_certificates.py
  paper/
    main.tex
    main.bbl
    source-map.tex
  release/
    manifest.sha256
\end{lstlisting}

Generator programs are intentionally outside the trusted proof path.  They may
be complex and use arbitrary precision, linear-programming, or interval
libraries.  Lean imports only certificate data and proves that the small
checkers accept it.

\subsection{Dependency order}

\begin{longtable}{p{0.065\textwidth}p{0.22\textwidth}p{0.365\textwidth}p{0.22\textwidth}}
\toprule
Layer & Product & Principal obligations & May depend on \\
\midrule
\endfirsthead
\toprule
Layer & Product & Principal obligations & May depend on \\
\midrule
\endhead
L0 & Definitions & Prime gaps, admissibility, support cells, weights, exact
parameters & Lean and mathlib only \\
L1 & Reusable foundations & Finite sums, arithmetic functions, congruences,
filters, integrals, exact polynomials & L0 and mathlib \\
L2 & Certificate checkers & Rational LP duality, rational interval and
polynomial checks, tuple check & L0--L1 \\
L3 & Finite witnesses & All 105 support pairs, trial function, \(I,J,K\),
minorant-loss bound, strict crossing & L0--L2 \\
L4 & Abstract sieve & Selberg/Maynard weights, \(S_1,S_2\), positivity, passage
from two primes in a tuple to a bounded consecutive gap & L0--L1 \\
L5 & Distribution & Bombieri--Vinogradov and the exact Type-I/II/III estimates,
including endpoints and uniformity & L0--L1 \\
L6 & Minorant & Harman minorant, roughness, mass, distribution on the verified
support & L1, L3, L5 \\
L7 & Assembly & Instantiate the abstract sieve and apply the tuple theorem &
L2--L6 \\
L8 & Audit & Axiom report, independent checker, paper/source map & all layers \\
\bottomrule
\end{longtable}

Keep imports one-way.  In particular, foundational analytic modules must not
import the final numerical data.  This prevents an accidental circular proof
or a theorem whose generality is only cosmetic.

\section{A staged proof strategy that exposes risk early}

Two strands should start immediately.  First, build a vertical formal skeleton
ending in a conditional theorem with an explicit \leanname{AnalyticInput}
argument.  Second, attack one representative hard analytic estimate early,
preferably the most demanding Type-III leaf identified by the source audit.
The first strand validates the finite architecture; the second prevents the
project from spending a year polishing certificates before discovering a
missing theory stack.

\subsection{Milestones and exit criteria}

\begin{longtable}{p{0.065\textwidth}p{0.20\textwidth}p{0.48\textwidth}p{0.125\textwidth}}
\toprule
ID & Milestone & Mandatory exit criterion & Status at start \\
\midrule
\endfirsthead
\toprule
ID & Milestone & Mandatory exit criterion & Status at start \\
\midrule
\endhead
M0 & Specification & Approved target statements, source ledger, pinned
toolchain, CI, and empty axiom audit & \todo \\
M1 & Conditional spine & A theorem
\leanname{conditional\_H1\_le\_236} whose only nonfinite hypothesis is one
precisely defined \leanname{AnalyticInput} structure & \todo \\
M2 & Finite support & Kernel-checked proofs of frontier reduction and all 105
partition pairs; no Python or LP solver in the trust boundary & \todo \\
M3 & Trial function & Exact formal witness and a proved strict inequality with
a rational margin; all integral enclosures checked in Lean & \todo \\
M4 & Abstract sieve & Maynard/Stadlmann implication proved from an explicit
distribution interface, with all limit and positivity steps & \todo \\
M5 & Classical distribution & Sufficient Bombieri--Vinogradov theorem proved,
including the precise weighted and dyadic form used later & \todo \\
M6 & Extended distribution & Type-I/II/III propositions proved with every
cited auxiliary estimate instantiated & \todo \\
M7 & Minorant & Constructed \(\rho\), bounds \(-24\leq\rho\leq1_{\mathbb P}\),
roughness, mass \(1-c_1\), and distribution & \todo \\
M8 & Unconditional theorem & \leanname{boundedPrimeGaps\_236} has no arguments,
builds cleanly, and passes the complete axiom audit & \todo \\
M9 & Release & Paper/source correspondence audit, independent rebuild,
archival code package, and arXiv preview approval & \todo \\
\bottomrule
\end{longtable}

M1 is deliberately conditional.  Define a structure instead of inserting
temporary axioms:

\begin{lstlisting}
structure AnalyticInput where
  rho : Nat -> Real
  lower_bound : ...
  upper_bound : ...
  rough_support : ...
  mass_asymptotic : ...
  equidistribution : ...

theorem conditional_H1_le_236
    (h : AnalyticInput) : BoundedPrimeGaps 236 := by
  ...

theorem constructAnalyticInput : AnalyticInput := by
  -- M5--M7.  This is the analytic mountain.
  ...

theorem boundedPrimeGaps_236 : BoundedPrimeGaps 236 :=
  conditional_H1_le_236 constructAnalyticInput
\end{lstlisting}

The ellipses shown in this guideline are specifications, not code permitted in
the repository.

\section{The analytic formalisation program}

\subsection{Source ledger before proof translation}

Every invocation of ``standard,'' ``similarly,'' or ``by a routine dyadic
decomposition'' in the sources must become an entry in
\leanname{SOURCE\_LEDGER.md}.  Each row should contain:
\[
\begin{gathered}
\text{paper locator},\quad \text{exact informal statement},\quad
\text{Lean theorem name},\\
\text{dependencies},\quad \text{endpoint conventions},\quad
\text{status}.
\end{gathered}
\]
If a citation gives a theorem in a different normalization, add a separate
conversion lemma.  Never hide normalization changes inside a long tactic
proof.

The attached Maynard paper is the published source for the basic small-gap
sieve \cite{Maynard2015}.  The existing \(236\) dossier then uses the
Stadlmann framework and a corrected negative-\(\omega_0\) branch
\cite{Stadlmann2026}.  Since the latter source is a manuscript, formalisation
must proceed from a corrected, self-contained mathematical statement rather
than from proposition numbers alone.

\subsection{Mathlib inventory}

At the pinned commit, use \leanname{\#check}, documentation search, and source
search to inventory at least:

\begin{itemize}
\item \leanname{Nat.Prime}, nth primes, prime counting, divisors, squarefree
integers, Möbius and von Mangoldt functions;
\item Dirichlet characters, finite Fourier sums, roots of unity, and character
orthogonality;
\item real and complex finite sums, Cauchy--Schwarz, Hölder, summation by parts,
and interval integrals;
\item filters and asymptotic notation at infinity;
\item existing Selberg-sieve structures and diagonalisation; and
\item any large-sieve, exponential-sum, Kloosterman-sum, or prime-distribution
theorems that can be reused without changing the desired statement.
\end{itemize}

Current documentation exhibits prime-counting, Dirichlet-character, Möbius,
asymptotic, and Selberg-sieve infrastructure
\cite{MathlibPrimeCounting,MathlibDirichlet,MathlibMoebius,MathlibAsymptotics,MathlibSelberg}.
It should not be assumed, without a pinned-source search, that a theorem named
Bombieri--Vinogradov or the needed Type-I/II/III estimates is available.

\subsection{Maynard sieve layer}

Formalise the narrowest theorem sufficient for \(k=48\), while keeping the
main combinatorial argument parameterized where this reduces duplication.  A
complete checklist is:

\begin{enumerate}[label=A\arabic*.,leftmargin=3.2em]
\item Define admissible tuples, the \(W\)-trick, residue selection, and the
  multidimensional divisor weights.
\item Define the sums \(S_1\) and \(S_2\), prove nonnegativity, and state the
  exact positivity implication that forces at least two prime values.
\item Expand squares and justify every interchange of finite sums.
\item Evaluate the main terms and isolate remainder terms with explicit
  uniform quantifiers.
\item Connect the distribution estimate for the minorant to the remainder
  bound needed by \(S_2\).
\item Prove that two primes in one translate of a diameter-\(236\) tuple contain
  a pair of consecutive primes at distance at most \(236\).  The two detected
  primes need not themselves be consecutive; choose adjacent primes inside
  their interval.
\item Convert arbitrarily large successful translates into
  \leanname{BoundedPrimeGaps 236}.
\end{enumerate}

The existing mathlib Selberg-sieve diagonalisation may shorten A2--A4, but its
upper-bound-sieve interface must be compared carefully with the lower-bound,
multidimensional weights used here.

\subsection{Bombieri--Vinogradov layer}

Prove only the version actually needed by the negative-\(\omega_0\) branch,
but do not weaken its statement by suppressing dependencies.  The theorem must
name:

\begin{itemize}
\item the interval or dyadic range for the integer variable;
\item the modulus range and the fixed positive distance below exponent
  \(1/2\);
\item whether weights are \(\Lambda\), \(\theta\), or prime indicators;
\item the coprimality and exceptional-residue conditions;
\item the norm in which errors are averaged over moduli;
\item the quantifier order for \(A\), the level loss, and the threshold
  \(x_0(A)\); and
\item the conversion from the proved form to the exact moving/dyadic form used
  in the sieve.
\end{itemize}

A conventional route is large sieve plus a combinatorial identity such as
Vaughan's identity.  Whether that is the shortest formal route must be decided
after the inventory.  Any quoted large-sieve inequality, zero-free result, or
character estimate not already proved in mathlib becomes a named subproject.

\subsection{Type-I/II/III layer}

This is likely the largest risk.  For each estimate:

\begin{enumerate}[label=D\arabic*.,leftmargin=3.2em]
\item State a finite-
\(x\) version with explicit constants or an eventual bound whose quantifiers
are completely formal.
\item Split the convolution variables exactly as in the corrected source.
\item Formalise the support-partition proposition separately from the analytic
estimate.  The certificate proves only that the exponents lie in an allowed
region; it does not prove the estimate on that region.
\item Trace every external exponential-sum input to a formal theorem.  If a
deep algebraic-geometric estimate is used only through a black-box paper
citation, it cannot be introduced as an axiom and still support the word
``unconditional.''
\item Prove uniformity under the finite collection of dyadic decompositions and
all smooth cutoffs.  Record powers of \(\log x\) instead of absorbing them
informally.
\item Close overlaps and endpoints explicitly, especially \(\omega_0=0\).
\end{enumerate}

Run an early feasibility spike on the hardest D4 input.  If its formal
dependency tree reaches an unavailable theory that would itself take years,
the project can either formalise that theory, seek a different analytic route,
or revise the target bound.  It must not silently declare the input.

\subsection{Harman minorant and mass}

Construct the minorant \(\rho(n;x)\) rather than postulating it.  The Lean
interface must prove:
\[
 -24\leq\rho(n;x)\leq\mathbf 1_{\mathbb P}(n),\qquad
 \rho(n;x)\ne0\Longrightarrow P^-(n)>x^{979/5000},
\]
the mass formula with loss \(c_1\), and precisely the distribution property
used by the sieve.  The rectangle majorant
\[
 c_1\leq
 \frac{1291475390625}{14389093733134384}
\]
is rational and should be proved directly from the exceptional-region
definitions.  In particular, prove rather than comment that every denominator
factor is at least \(979/5000\).

\section{Formalising the finite support certificate}

The current certificate uses two support slices, \(105\) frontier pairs,
\(73\) direct Type-IIc cases, \(8\) adaptive cases, and \(24\) rational-LP
cases.  These counts should appear as Lean theorems, not comments.

\subsection{Exact parameters}

Represent every terminating decimal as a rational literal or a numerator and
denominator.  Define the parameter bundle once:
\[
 k=48,\quad \delta=\frac7{500},\quad
 A_1=\frac{253}{1000},\quad A_2=\frac{129}{500},
\]
\[
 \xi_1=\frac{759}{2000},\qquad
 \xi_2=\xi_3=\frac{4021}{10000},\qquad
 \eta_P=10^{-10}.
\]
Do not let real-number coercions obscure an exact inequality: first prove the
rational comparison, then transport it to \(\mathbb R\).

\subsection{Frontier reduction}

Prove the merge lemma abstractly: if a block has more than six entries, merge
two entries; the merged mass remains at most \(0.202<1\), and any packing of the
merged list lifts to a packing of the original list by keeping the two
constituents in the same bin.  Iteration reduces the problem to counts
\(0,\ldots,6\), producing
\[
 \binom82+7^2+\binom82=105
\]
unordered lower/lower, lower/upper, and upper/upper pairs.

Model bins and assignments with finite types.  Separate the mathematical
predicate \leanname{Fits} from the decision procedure \leanname{checkFits}, and
prove a soundness theorem
\[
 \leanname{checkFits cert input = true}\Longrightarrow
 \leanname{Fits input}.
\]

\subsection{LP dual checker}

An external LP solver may propose partitions and dual multipliers.  The Lean
checker must verify over \(\mathbb Q\):

\begin{enumerate}
\item dimensions and index bounds;
\item signs of all dual multipliers;
\item exact stationarity \(A^{T}y=c\);
\item exact dual objective comparison; and
\item the conclusion by a proved weak-duality lemma.
\end{enumerate}

Certificate deserialization is not a logical step.  Prefer generated Lean data
files whose arrays have statically checked lengths, or parse a simple format
and prove that checker acceptance implies the theorem regardless of parser
behavior.  Floating-point output is never admitted directly.

\subsection{Axiom-free computation policy}

For modest computations, prefer definitional reduction,
\leanname{norm\_num}, or kernel-evaluated decision procedures.  Native
evaluation can enlarge the trust base and is visible in the axiom report; the
exact behavior is Lean-version dependent \cite{LeanValidation}.  The release
policy is therefore:

\begin{center}
\begin{tabularx}{\textwidth}{p{0.29\textwidth}X p{0.18\textwidth}}
\toprule
Method & Policy & Release status \\
\midrule
Explicit proof term / \leanname{norm\_num} & Preferred where practical &
\pass \\
Kernel reduction / proved reflective checker & Preferred for certificates &
\pass \\
Native computation producing a reported computation axiom & Useful for
exploration, replace before strict release & \fail \\
External Arb/LP assertion imported as a theorem & Never acceptable & \fail \\
Project \leanname{axiom}, \leanname{sorry}, or \leanname{admit} & Never
acceptable & \fail \\
\bottomrule
\end{tabularx}
\end{center}

If kernel evaluation is too slow, improve the certificate format and checker;
do not weaken the theorem's trust claim.

\section{Formalising the trial function and strict crossing}

\subsection{Exact Chebyshev witness}

Define the three piecewise polynomial profiles from the archived rational
Chebyshev coefficients.  Prove the Chebyshev recurrence and the equivalence of
the Chebyshev and monomial evaluations used by the certificate.  Define the
symmetric trial function from these profiles and prove measurability, bounded
support, symmetry, and square integrability.

The current proof gives an exact rational \(I\)-quadratic form and tight
outward-rounded enclosures for the \(J\)-matrix.  In Lean, replace the Arb
trust assumption by one of these proof-producing routes, in preferred order:

\begin{enumerate}
\item exact rational integration of polynomial pieces over rational cells;
\item rational Bernstein bounds on a certified subdivision, with a proved
  integration bound; or
\item a small rational interval-arithmetic checker proved sound in Lean.
\end{enumerate}

The generator may use Arb to choose subdivisions, but the final lower and upper
bounds must follow from rational certificates checked by Lean.  Keep the strict
margin as a rational inequality after clearing positive denominators.

\subsection{Prefer an explicit positive enlargement}

The present dossier proves the crossing first at support-enlargement parameter
\(\varepsilon=0\), then invokes dominated convergence and absolute continuity
to obtain some \(\varepsilon>0\).  That is valid as an informal route but
expensive to formalise and easy to mismatch with open/closed boundaries.

The preferred formal route is to choose a concrete positive rational
\(\varepsilon_*\), small enough to retain the margin, and certify directly:
\[
 \frac{48(1-c_1)J(F;\varepsilon_*)-
       48\cdot24 K(F;\varepsilon_*)}{I(F;\varepsilon_*)}>1.
\]
Do not guess \(\varepsilon_*\).  Extend the existing exact geometry program to
search for one, then import only the rational certificate into Lean.  This
removes an entire measure-theoretic limit bridge from the critical path.  If no
reasonable direct certificate is found, formalise the convergence proof with
all boundary-null-set and domination lemmas.

\begin{gate}[Finite crossing]
The gate passes only when one theorem with no hypotheses proves the exact
strict sieve inequality for a specific rational witness and, preferably, a
specific positive rational \(\varepsilon_*\).  Decimal display values are
documentation only.
\end{gate}

\section{Controlling asymptotics and quantifiers}

Informal analytic-number-theory papers routinely reuse \(o(1)\), allow implied
constants to change, and choose thresholds in a nested order.  Lean forces the
order to be correct.  Maintain a parameter ledger such as
\[
 A\longmapsto \delta_A\longmapsto x_0(A,\delta_A),
\]
and never reuse one existential threshold for a stronger statement without a
maximum argument.

Prefer eventual finite inequalities over a proliferation of bespoke
little-\(o\) definitions.  For example, the final positivity step should have
the shape
\[
 \exists x_0\ \forall x\geq x_0,
 \quad S_2(x)-S_1(x)>0,
\]
deduced from named main-term bounds.  Generic filter theorems can establish the
eventual bounds, but the final arithmetic implication should be explicit.

For each use of asymptotic notation, record:
\begin{itemize}
\item domain and filter;
\item all fixed and varying parameters;
\item positivity conditions for logarithms and denominators;
\item uniformity over residues, moduli, tuple indices, and support cells; and
\item the exact later inequality for which the estimate is sufficient.
\end{itemize}

\section{ChatGPT/Codex protocol inside VS Code}

AI assistance is most reliable when each request has one mathematical target,
one file boundary, and an executable acceptance test.  Do not ask an agent to
``formalise Section 6'' in one pass.

\subsection{Required task packet}

Every Codex task should include:

\begin{enumerate}
\item exact file and theorem name;
\item current imports and the pinned toolchain;
\item the fully quantified mathematical statement;
\item the relevant source paragraph or a self-contained proof;
\item permitted helper lemmas and prohibited trust shortcuts;
\item the command that must pass; and
\item the expected axiom audit.
\end{enumerate}

Use the following prompt template.

\begin{lstlisting}[language={}]
You are working in the pinned Lean 4/mathlib project PrimeGaps236.

Task: prove theorem <fully-qualified-name> in <path>.
Mathematical statement: <fully quantified statement>.
Informal proof/source: <self-contained proof or exact locator>.
Allowed dependencies: <list>.

Requirements:
1. Inspect existing declarations with #check and source search; do not invent APIs.
2. Do not use axiom, sorry, admit, unsafe theorem insertion, or native assertions.
3. Preserve the intended quantifier order and endpoint conventions.
4. Add small named helper lemmas when they expose mathematical structure.
5. Run: lake build <Module.Name>
6. Run the relevant #print axioms audit and report its exact output.
7. Summarize any mismatch between the informal proof and the formal statement.

Do not weaken the theorem or add hypotheses without explaining the blocker.
\end{lstlisting}

\subsection{Four useful agent modes}

\begin{longtable}{p{0.18\textwidth}p{0.32\textwidth}p{0.40\textwidth}}
\toprule
Mode & Request & Required result \\
\midrule
Implementation & Prove one approved lemma & Compiling change plus local axiom
audit \\
API reconnaissance & Find existing mathlib support for one concept & Exact
declaration names, imports, and a compiling scratch example; no project edit \\
Adversarial audit & Try to falsify statement or locate hidden assumptions &
Concrete counterexample, missing hypothesis, or checked confirmation \\
Refactor & Shorten or generalize already proved code & Identical public theorem
and unchanged axiom output \\
\bottomrule
\end{longtable}

Separate implementation and adversarial-audit chats.  The auditor should be
given the statement and source but not encouraged to preserve the implementer's
approach.  For high-risk analytic lemmas, a human must review both the Lean
statement and the mathematical reduction; a green build checks the term, not
the intended interpretation.

\subsection{When Codex gets stuck}

After two or three failed proof shapes, stop local tactic experimentation and
classify the obstacle:

\begin{enumerate}
\item missing library theorem;
\item wrong or under-specified mathematical statement;
\item coercion/normalization mismatch;
\item inadequate data representation;
\item performance problem in reduction; or
\item a genuine gap in the paper proof.
\end{enumerate}

For (1), prove a reusable lemma or consult the mathlib community.  For (2) or
(6), return to the source ledger and obtain mathematical review.  For (4) or
(5), redesign the interface or certificate.  Repeated tactic search should not
be used to conceal an unresolved mathematical premise.

\section{Continuous verification and trust audit}

\subsection{Local checks}

Run at least the following before every merge:

\begin{lstlisting}[language=bash]
lake build
rg -n '\bsorry\b|\badmit\b|^\s*axiom\b' PrimeGaps236
git diff --check
\end{lstlisting}

Keep an audit module containing:

\begin{lstlisting}
import PrimeGaps236.Final

#check PrimeGaps236.boundedPrimeGaps_236
#print axioms PrimeGaps236.boundedPrimeGaps_236
\end{lstlisting}

The exact output belongs in the release log.  A text search is only a first
line of defense; the transitive axiom report is authoritative for accidental
imports of incomplete results.

\subsection{Independent rechecking}

Lean's validation guide recommends \leanname{lean4checker --fresh} on the
built module and identifies comparator plus an external checker as the gold
standard for unreviewed or potentially adversarial generated code
\cite{LeanValidation}.  Accordingly, the release process should include:

\begin{enumerate}
\item a clean build on two independently provisioned machines;
\item \leanname{lean4checker --fresh} for the final module;
\item a sandboxed comparator/external-checker run if supported by the pinned
Lean version;
\item a manual raw inspection of the final theorem statement and the unfolded
definitions it uses; and
\item SHA-256 hashes for the source bundle and verification logs.
\end{enumerate}

Never run unreviewed generator or tactic code with valuable credentials in the
environment.  Proof-generating metaprograms are executable programs even when
their output is later kernel checked.

\begin{gate}[Unconditional theorem]
The gate passes only if a clean checkout builds
\leanname{PrimeGaps236.Final}, the primary theorem has no arguments, the axiom
output contains no nonstandard/project/computation axiom, and an independent
review confirms that the formal statement means infinitely many consecutive
prime gaps of size at most \(236\).
\end{gate}

\section{Paper-to-code correspondence}

The mathematical paper should be readable without Lean, and the Lean code
should be checkable without trusting prose.  Link them with a source map.

For every numbered theorem or nontrivial computational claim, the paper should
give:
\begin{itemize}
\item its fully qualified Lean theorem name;
\item the source module;
\item whether the proof is conceptual, reflective, or certificate-based; and
\item the relevant certificate identifier and checksum, if any.
\end{itemize}

A generated appendix table may have columns:

\begin{center}
\small
\begin{tabularx}{\textwidth}{>{\raggedright\arraybackslash}p{0.15\textwidth}>{\raggedright\arraybackslash}p{0.27\textwidth}>{\raggedright\arraybackslash}p{0.23\textwidth}>{\raggedright\arraybackslash}X}
\toprule
Paper item & Lean declaration & Module & Dependencies \\
\midrule
Main theorem & \leanname{boundedPrimeGaps\_236} &
\leanname{PrimeGaps236.Final} & Analytic input construction, crossing, tuple \\
Sieve criterion & \leanname{maynardCriterion} &
\path{Sieve.MaynardCriterion} & Main-term and distribution lemmas \\
Support certificate & \leanname{allSupportCases} &
\path{Certificates.SupportData} & Sound rational checker \\
Numerical crossing & \leanname{trialCrossing} &
\path{TrialFunction.Crossing} & Exact rational integral certificates \\
Tuple & \leanname{admissibleTuple236} &
\path{Tuple.Admissible236} & Finite residue computation \\
\bottomrule
\end{tabularx}
\end{center}

The paper's abstract must distinguish three dates: the date the mathematical
proof was completed, the Lean release tag, and the arXiv version.  State the
trusted computing base accurately.  Do not say ``verified by Lean'' if the
final theorem still takes \leanname{AnalyticInput} as an argument or if native
certificate axioms appear in its dependency closure.

\section{arXiv release package}

arXiv supports program code as ancillary material only when the article is
submitted from TeX/PDFLaTeX source, not as a PDF-only submission.  Ancillary
files belong under an \leanname{anc/} directory at the root of the upload
package \cite{ArxivAncillary}.  Use this structure:

\begin{lstlisting}[language={}]
arxiv-upload/
  main.tex
  main.bbl
  figures/
  00README
  anc/
    README.txt
    lean/
      lean-toolchain
      lakefile.toml
      lake-manifest.json
      PrimeGaps236.lean
      PrimeGaps236/
        ... all imported .lean source files ...
      REPRODUCING.md
      LICENSE
      manifest.sha256
\end{lstlisting}

Do not put paper TeX files in \leanname{anc/}.  Do not rely on links from code
that contain the literal \leanname{anc/} prefix, because arXiv may change the
download layout.  Do not include \leanname{.git/}, \leanname{.lake/},
\leanname{.vscode/}, generated \leanname{.olean} files, local caches, solver
binaries, or secrets.  arXiv deletes hidden files on announcement, so the
ancillary project must reproduce without them \cite{ArxivTeX}.

The root \leanname{00README} should identify \leanname{main.tex} as the top
level file if necessary.  Compile with an arXiv-supported TeX Live setup,
include the matching \leanname{main.bbl}, omit auxiliary build products, and
inspect the PDF that arXiv itself generates.  arXiv currently advises against
extraneous source files and requires authors to review the processed PDF
\cite{ArxivTeX}.

The ancillary \leanname{README.txt} must contain exact commands:

\begin{lstlisting}[language={}]
1. Install elan, git, and curl.
2. In this directory run: lake exe cache get
3. Run: lake build
4. Run the final axiom audit documented in AxiomAudit.lean.
5. Run lean4checker --fresh on PrimeGaps236.Final if available.
Expected theorem: PrimeGaps236.boundedPrimeGaps_236
Expected statement: BoundedPrimeGaps 236
\end{lstlisting}

Also create a public version-controlled repository and archive its exact
release tag with a durable identifier.  The arXiv ancillary copy remains the
version-specific snapshot; the repository supports issue tracking and later
toolchain ports.  The paper should link both and state their hashes.

\section{Release checklist}

No item may be waived silently.

\begin{longtable}{p{0.05\textwidth}p{0.73\textwidth}p{0.12\textwidth}}
\toprule
ID & Check & Required \\
\midrule
\endfirsthead
\toprule
ID & Check & Required \\
\midrule
\endhead
R1 & The unfolded final statement is the intended assertion about infinitely
many consecutive prime gaps. & \mustpass \\
R2 & The final theorem has no analytic hypothesis or certificate hypothesis. &
\mustpass \\
R3 & Every source citation used mathematically maps to a proved Lean theorem. &
\mustpass \\
R4 & The negative-\(\omega_0\) branch and the \(\omega_0=0\) endpoint are
proved in their exact modulus normalization. & \mustpass \\
R5 & The 105 support cases and their \(73+8+24\) split are checked by proved
rational procedures. & \mustpass \\
R6 & The trial-function crossing is an exact rational theorem; no external
ball or LP output is trusted. & \mustpass \\
R7 & The explicit 48-tuple is formally admissible and has diameter 236. &
\mustpass \\
R8 & \leanname{lake build} succeeds in a clean checkout from the archived
files. & \mustpass \\
R9 & \leanname{\#print axioms} contains only the approved standard foundational
axioms. & \mustpass \\
R10 & \leanname{lean4checker --fresh} and, if available, external checking
succeed. & \mustpass \\
R11 & Two mathematical reviewers audit the analytic statements and one
independent Lean reviewer audits the trust boundary. & \mustpass \\
R12 & Paper theorem numbers, Lean names, release tag, and hashes agree. &
\mustpass \\
R13 & The arXiv source compiles on the current arXiv TeX system and the generated
PDF is visually inspected. & \mustpass \\
R14 & Ancillary code contains all sources and manifests, but no hidden caches,
binaries, credentials, or paper TeX files. & \mustpass \\
R15 & The title and abstract accurately say ``unconditional Lean
formalisation'' only after R1--R14 pass. & \mustpass \\
\bottomrule
\end{longtable}

\section{Gap register and decision rules}

Maintain a public \leanname{GAPS.md}.  At minimum begin with:

\begin{longtable}{p{0.06\textwidth}p{0.275\textwidth}p{0.375\textwidth}p{0.16\textwidth}}
\toprule
ID & Risk & Required resolution & Priority \\
\midrule
\endfirsthead
\toprule
ID & Risk & Required resolution & Priority \\
\midrule
\endhead
G1 & Stadlmann analytic framework is not presently a proved Lean dependency &
Formalise every proposition and cited input in the corrected normalization &
Critical \\
G2 & Printed negative-\(\omega_0\) Type-IIc capacity is impossible & Prove
the separate Bombieri--Vinogradov branch and the endpoint split & Critical \\
G3 & Deep Type-III inputs may require absent foundational theory & Complete an
early vertical feasibility proof and expand the source ledger & Critical \\
G4 & Existing support LP logs do not themselves constitute Lean proofs & Export
full rational witnesses and verify with a proved checker & High \\
G5 & Existing \(J\)-bounds trust rational geometry plus Arb & Replace by exact
rational or proved interval certificates & High \\
G6 & Passage from \(\varepsilon=0\) to positive \(\varepsilon\) is non-effective
& Prefer a directly certified rational \(\varepsilon_* >0\), or formalise the
limit argument & High \\
G7 & Informal \(o(1)\) terms conceal quantifier order & State uniform eventual
bounds and prove the final positivity threshold & High \\
G8 & AI-generated statements may be subtly weaker than the paper & Independent
statement review, raw unfolding, and paper-to-code map & Critical \\
\bottomrule
\end{longtable}

Use these decision rules:

\begin{enumerate}
\item A discovered mathematical gap pauses downstream theorem claims but does
not justify an axiom.
\item A slow certificate triggers representation work before more hardware.
\item An API mismatch triggers a small adapter lemma, not a rewrite of the
mathematical statement.
\item A source estimate with unavailable prerequisites triggers a documented
go/no-go review after the early feasibility spike.
\item The numerical target may be weakened only by an explicit project
decision and a new statement review; Codex must not weaken it to make a proof
compile.
\end{enumerate}

\section{Recommended immediate next sprint}

The first sprint should produce evidence about feasibility rather than a large
volume of code.

\begin{enumerate}[label=S\arabic*.,leftmargin=3.2em]
\item Create the pinned repository and implement the direct
\leanname{BoundedPrimeGaps} statement plus its nth-prime bridge.
\item Translate the explicit tuple and prove its admissibility without native
trust axioms.
\item Write \leanname{AnalyticInput} with no omitted quantifiers, then compile a
conditional final spine whose finite requirements are still named structure
fields.
\item Complete the source ledger for the negative-\(\omega_0\) branch and one
hard Type-III estimate, recursively listing every cited theorem.
\item Implement a tiny prototype rational LP checker and verify one of the 24
saved dual cases end to end.
\item Search for and certify a concrete positive \(\varepsilon_*\), because
this can remove the limiting-support bridge.
\item At sprint end, publish a dependency report: proved, reusable from
mathlib, needs formalisation, or mathematically unresolved.
\end{enumerate}

This sprint should end with a realistic estimate of the analytic mountain.  It
should not end with a public claim that Lean has proved \(\Hone\leq236\).

\section{Effort model}

Calendar predictions are unreliable before the source audit, but the relative
scale is clear.

\begin{center}
\begin{tabularx}{\textwidth}{p{0.31\textwidth}p{0.18\textwidth}X}
\toprule
Work package & Relative size & Main uncertainty \\
\midrule
Statement, tuple, repository & Small & nth-prime bridge and definitions \\
Support and LP certificates & Medium & certificate volume and reduction speed \\
Trial function and integrals & Medium--large & efficient exact geometry checker \\
Abstract Maynard sieve & Large & asymptotic main-term infrastructure \\
Bombieri--Vinogradov & Large & large-sieve and character-sum prerequisites \\
Harman minorant & Large & combinatorial identities and uniform mass bounds \\
Type-I/II/III distribution & Very large & deep exponential-sum dependencies \\
Audit, paper, and release & Medium & independent review and reproducibility \\
\bottomrule
\end{tabularx}
\end{center}

A small expert team should expect months for the finite and abstract-sieve
layers and potentially years for the complete analytic chain.  This is not a
reason to avoid the project; it is a reason to perform the hardest dependency
spike early and to describe intermediate releases accurately.

\section{Definition of done}

The project is complete exactly when the archived sources can be placed in a
fresh environment and the following chain succeeds:
\[
\begin{gathered}
\text{pinned toolchain}\ \longrightarrow\ \text{clean build}
\longrightarrow\ \text{final no-argument theorem},\\
\text{approved axiom report}\ \longrightarrow\ \text{independent replay},\\[2pt]
\texttt{PrimeGaps236.boundedPrimeGaps\_236}:
\texttt{BoundedPrimeGaps 236}.
\end{gathered}
\]
The paper must then explain that theorem, the corrected analytic proof, and the
certificate checkers in enough mathematical detail for independent experts to
understand what Lean checked.  The arXiv ancillary source must reproduce the
same theorem at the same release tag.  Anything short of this is a valuable
intermediate formalisation, but not the requested end-to-end unconditional
formalisation.

\begin{thebibliography}{99}

\bibitem{Maynard2015}
J. Maynard,
\emph{Small gaps between primes},
Annals of Mathematics \textbf{181} (2015), 383--413.
\url{https://doi.org/10.4007/annals.2015.181.1.7}.

\bibitem{Stadlmann2026}
J. Stadlmann,
\emph{On the smallest gaps between primes},
arXiv:2608.31126 (manuscript version to be fixed in the project source ledger).
\url{https://arxiv.org/abs/2608.31126}.

\bibitem{LeanInstall}
Lean project,
\emph{Install Lean and create a mathlib project},
accessed 3 September 2026.
\url{https://lean-lang.org/install/manual/}.

\bibitem{LeanValidation}
Lean project,
\emph{Validating a Lean Proof},
Lean Language Reference, accessed 3 September 2026.
\url{https://lean-lang.org/doc/reference/latest/ValidatingProofs/}.

\bibitem{MathlibReadme}
Lean community,
\emph{mathlib4: installation, documentation, and build commands},
accessed 3 September 2026.
\url{https://github.com/leanprover-community/mathlib4/blob/master/README.md}.

\bibitem{MathlibPrimeCounting}
Lean community,
\emph{Mathlib.NumberTheory.PrimeCounting},
mathlib documentation, accessed 3 September 2026.
\url{https://leanprover-community.github.io/mathlib4_docs/Mathlib/NumberTheory/PrimeCounting.html}.

\bibitem{MathlibSelberg}
Lean community,
\emph{Mathlib.NumberTheory.SelbergSieve},
mathlib documentation, accessed 3 September 2026.
\url{https://leanprover-community.github.io/mathlib4_docs/Mathlib/NumberTheory/SelbergSieve.html}.

\bibitem{MathlibDirichlet}
Lean community,
\emph{Mathlib.NumberTheory.DirichletCharacter.Basic},
mathlib documentation, accessed 3 September 2026.
\url{https://leanprover-community.github.io/mathlib4_docs/Mathlib/NumberTheory/DirichletCharacter/Basic.html}.

\bibitem{MathlibMoebius}
Lean community,
\emph{Mathlib.NumberTheory.ArithmeticFunction.Moebius},
mathlib documentation, accessed 3 September 2026.
\url{https://leanprover-community.github.io/mathlib4_docs/Mathlib/NumberTheory/ArithmeticFunction/Moebius.html}.

\bibitem{MathlibAsymptotics}
Lean community,
\emph{Mathlib.Analysis.Asymptotics.Basic},
mathlib documentation, accessed 3 September 2026.
\url{https://leanprover-community.github.io/mathlib4_docs/Mathlib/Analysis/Asymptotics/Basic.html}.

\bibitem{MathInLean}
J. Avigad, P. Massot, F. van Doorn, and others,
\emph{Mathematics in Lean}, accessed 3 September 2026.
\url{https://leanprover-community.github.io/mathematics_in_lean/}.

\bibitem{ArxivAncillary}
arXiv,
\emph{Ancillary Files (data, code, images)},
accessed 3 September 2026.
\url{https://info.arxiv.org/help/ancillary_files.html}.

\bibitem{ArxivTeX}
arXiv,
\emph{Submit TeX/LaTeX},
accessed 3 September 2026.
\url{https://info.arxiv.org/help/submit_tex.html}.

\end{thebibliography}

\end{document}
